\section{Introduction}\label{sec:intro}

Many questions asked of a large linear dynamical system concern a handful of its outputs. A thermal model of a chip is probed at a few sensor locations while a few heaters are driven; a networked system is excited at a few driver nodes and observed at a few others~\cite{liu2011,pasqualetti2014,mesbahi2010}; a stochastic model is queried for the probability of a few configurations. In state-space form the object of interest is $\by(t)=C\bx(t)$ for $\dot\bx=A\bx+B\bu$, with $C$ and $B$ having few rows and columns, and the state dimension $N$ may be far larger than anything that has to be reported.

Standard methods nevertheless work with the full state. Time integrators propagate $\bx(t)$. The scaled Taylor method of Al-Mohy and Higham~\cite{almohy2011} and Krylov methods~\cite{hochbruck1997,sidje1998} compute the action $e^{tA}\bv$ and need a few, or $m+1$, vectors of length $N$, in addition to $A$ itself; limited-memory variants reduce the multiple but not the scaling~\cite{guttel2020}. None of these stores $e^{tA}$, and for moderate $N$ they are the right tools. They become the wrong tools when the state vector itself does not fit, which happens for implicitly specified models with astronomically many states, or when memory is the binding constraint and only a few numbers are wanted. Local methods restrict the computation to a neighborhood of the inputs, as in the finite state projection for master equations~\cite{munsky2006} or push methods for exponentials of graph matrices~\cite{kloster2014,gleich2015}, and their memory is the size of that neighborhood. Model reduction~\cite{antoulas2005} amortizes an offline cost over many queries and is a different trade.

Barash, G\"uttel and Hen~\cite{barash2022} showed that an individual entry of a matrix function can be evaluated as a sum over walks on the graph of off-diagonal couplings, each walk weighted by a divided difference of the function at the diagonal entries it visits, with memory that does not grow with $N$. This is the off-diagonal (permutation-matrix) expansion developed for quantum Monte Carlo~\cite{albash2017,gupta2020pmr}; related path-sum constructions appear in Refs.~\cite{giscard2013,giscard2015,gorban2011}. Here we develop this representation for linear time-invariant systems, homogeneous and forced, and ask a question with a measurable answer: when can selected input--output responses be computed from walk sums with working memory independent of $N$, with a certified truncation error and at a practical cost, and how does that compare with the alternatives, in particular with local methods that use the same error bounds?

What is established and what is new should be kept apart. The walk expansion of $e^{tA}$ and its depth-first evaluation with incremental divided differences are those of Ref.~\cite{barash2022}. That an exponential forcing can be absorbed by enlarging the system is classical~\cite{vanloan1978,almohy2011}. The contributions of this paper are the following. (i) A linear input--output calculus in walk form: responses to exponential-polynomial inputs append a block of confluent nodes to each walk (Proposition~\ref{prop:forcing}), delayed and piecewise inputs are handled through a truncated-power basis without reconstructing the state at breakpoints, and all input patterns and observation times are obtained from one traversal. (ii) An output-oriented algorithm: backward enumeration from the sensors, lazy evaluation of the temporal factors, and a divided-difference kernel that returns a whole time grid by exact stepping of each walk's bidiagonal system; it handles repeated and clustered nodes without special cases, uses only nonnegative quantities in its squaring and stepping stages for real nodes, and was tested for complex ones. (iii) Error control: uniform-in-time componentwise bounds in terms of the comparison matrix $H^{-1}|R|$ that relates couplings to local decay rates, finite-horizon bounds for the general case, and an a posteriori certificate of the truncation error that accompanies subtree pruning, together with an accounting of the floating-point errors that the certificate does not cover. (iv) The same bounds certify local truncations of the system, a ball around the inputs and an adaptively grown set (Proposition~\ref{prop:local}), which gives local comparators with the same guarantee. (v) A measured assessment of memory and time from scratch against global and local methods on networks of up to $2^{64}$ states, with shared distance information, including the regimes where the walk sum loses. Appendix~\ref{app:td} records the extension to generators with exponential time dependence.

The measurements support three statements. First, the method does what its memory accounting promises: a few megabytes for the whole process, flat in $N$, for responses with a certified truncation error of systems whose state could not be stored. Second, on locally tree-like coupling graphs of degree 16 and 32 it was also the fastest of the certified methods we tested, against the uniform ball and against the adaptive set, given the same distance information, at every target that any of them reached; its peak memory stayed between 4.5 and 6.1~MB (7.2~MB with the larger distance map used for separated sensors) while the local methods needed up to 3~GB. The advantage held over three graph seeds and three actuator locations per seed for each degree, by factors of 1.5 to 19 in time against the adaptive set. At degree 4 all methods finished in tens of milliseconds and the adaptive set was marginally faster. Third, the walk sum loses, by orders of magnitude, on lattices, where many walks revisit the same states and a local state vector merges them, and when the coupling-to-decay ratio approaches or exceeds one, as in strongly nonnormal transport. Section~\ref{sec:numerics} documents these statements. In short, for selected responses of implicitly represented systems, certified walk pruning replaces the retained state vector of a local method by a small traversal workspace, and on high-degree tree-like coupling graphs the lazy evaluation of temporal factors makes it faster as well; the $2^{64}$-state examples illustrate the access model rather than the source of the advantage.
